\section{Projective covers, radicals, and finiteness conditions}

\begin{corollary}
    For finitely generated projective modules $P$ and $Q$,
    \[
        P\cong Q
        \Longleftrightarrow
        \operatorname{Top}(P)\cong \operatorname{Top}(Q).
    \]
\end{corollary}

\begin{proof}
    Let
    \[
        T := P/\operatorname{Rad}(P).
    \]
    By Nakayama lemma, the quotient map $P \to T$ is an essential epimorphism, so $P$ is the projective cover of $T$.
    Similarly $Q \to Q/\operatorname{Rad}(Q)$ is a projective cover.
    Then apply uniqueness of projective covers from the previous lecture.
\end{proof}

\begin{definition}
    A unital ring $R$ is called left Noetherian if ${}_R R$ is Noetherian as a left $R$-module.
\end{definition}

\begin{remark}
    Similar definitions are given for right Noetherian and for left/right Artinian rings.
\end{remark}

\begin{remark}
    one-sided Noetherian / Artinian conditions generally do not imply the opposite-sided versions
\end{remark}

\begin{example}
    There exists a ring which is left Noetherian but not right Noetherian.

    Let
    \[
        R=
        \left\{
        \begin{pmatrix}
            a & b \\
            0 & n
        \end{pmatrix}
        \;\middle|\;
        a,b\in \mathbb{Q},\ n\in \mathbb{Z}
        \right\}.
    \]
    Then $R$ is left Noetherian, but not right Noetherian.
\end{example}

\begin{proof}
    Let
    \[
        e_{11}=
        \begin{pmatrix}
            1 & 0 \\
            0 & 0
        \end{pmatrix},
        \qquad
        e_{22}=
        \begin{pmatrix}
            0 & 0 \\
            0 & 1
        \end{pmatrix}.
    \]
    Then, as left $R$-modules,
    \[
        {}_R R = Re_{11}\oplus Re_{22}.
    \]

    We first show that both $Re_{11}$ and $Re_{22}$ are Noetherian as left
    $R$-modules.

    For $Re_{11}$, we have
    \[
        Re_{11}
        =
        \left\{
        \begin{pmatrix}
            q & 0 \\
            0 & 0
        \end{pmatrix}
        \;\middle|\;
        q\in\mathbb{Q}
        \right\}.
    \]
    This is a simple left $R$-module. Indeed, if
    \[
        0\neq
        \begin{pmatrix}
            q & 0 \\
            0 & 0
        \end{pmatrix}
        \in Re_{11},
    \]
    then for any $x\in \mathbb{Q}$,
    \[
        \begin{pmatrix}
            xq^{-1} & 0 \\
            0       & 0
        \end{pmatrix}
        \begin{pmatrix}
            q & 0 \\
            0 & 0
        \end{pmatrix}
        =
        \begin{pmatrix}
            x & 0 \\
            0 & 0
        \end{pmatrix}.
    \]
    Hence every nonzero submodule of $Re_{11}$ is all of $Re_{11}$, so
    $Re_{11}$ is simple, in particular Noetherian.

    Now consider
    \[
        Re_{22}
        =
        \left\{
        \begin{pmatrix}
            0 & q \\
            0 & n
        \end{pmatrix}
        \;\middle|\;
        q\in\mathbb{Q},\ n\in\mathbb{Z}
        \right\}.
    \]
    Define
    \[
        \pi:Re_{22}\longrightarrow \mathbb{Z},
        \qquad
        \pi\!\left(
        \begin{pmatrix}
                0 & q \\
                0 & n
            \end{pmatrix}
        \right)=n.
    \]
    We regard $\mathbb{Z}$ as a left $R$-module via
    \[
        \begin{pmatrix}
            a & b \\
            0 & n
        \end{pmatrix}\cdot m = nm.
    \]
    Then $\pi$ is a surjective $R$-homomorphism. Its kernel is
    \[
        \ker\pi=
        \left\{
        \begin{pmatrix}
            0 & q \\
            0 & 0
        \end{pmatrix}
        \;\middle|\;
        q\in\mathbb{Q}
        \right\}.
    \]
    This is again a simple left $R$-module, because
    \[
        \begin{pmatrix}
            a & b \\
            0 & n
        \end{pmatrix}
        \begin{pmatrix}
            0 & q \\
            0 & 0
        \end{pmatrix}
        =
        \begin{pmatrix}
            0 & aq \\
            0 & 0
        \end{pmatrix},
    \]
    so the action is through multiplication by elements of $\mathbb{Q}$.
    Moreover, ${}_R\mathbb{Z}$ is Noetherian, since its submodules are exactly
    the ideals of $\mathbb{Z}$, and $\mathbb{Z}$ is Noetherian.

    Therefore $Re_{22}$ is an extension of two Noetherian left $R$-modules,
    hence is Noetherian. It follows that
    \[
        {}_R R = Re_{11}\oplus Re_{22}
    \]
    is a Noetherian left $R$-module. Thus $R$ is left Noetherian.

    Next we show that $R$ is not right Noetherian.

    For each integer $k\ge 0$, define
    \[
        I_k=
        \left\{
        \begin{pmatrix}
            0 & q \\
            0 & 0
        \end{pmatrix}
        \;\middle|\;
        q\in 2^{-k}\mathbb{Z}
        \right\}.
    \]
    We claim that each $I_k$ is a right ideal of $R$. Indeed, if
    \[
        \begin{pmatrix}
            0 & q \\
            0 & 0
        \end{pmatrix}\in I_k,
        \qquad
        \begin{pmatrix}
            a & b \\
            0 & n
        \end{pmatrix}\in R,
    \]
    then
    \[
        \begin{pmatrix}
            0 & q \\
            0 & 0
        \end{pmatrix}
        \begin{pmatrix}
            a & b \\
            0 & n
        \end{pmatrix}
        =
        \begin{pmatrix}
            0 & qn \\
            0 & 0
        \end{pmatrix}.
    \]
    Since $q\in 2^{-k}\mathbb{Z}$ and $n\in\mathbb{Z}$, we have
    $qn\in 2^{-k}\mathbb{Z}$, so the product again lies in $I_k$.

    Moreover,
    \[
        I_0\subsetneq I_1\subsetneq I_2\subsetneq \cdots
    \]
    is a strictly ascending chain of right ideals, because
    \[
        2^{-(k+1)}\in 2^{-(k+1)}\mathbb{Z}
        \quad\text{but}\quad
        2^{-(k+1)}\notin 2^{-k}\mathbb{Z}.
    \]
    Hence $R$ does not satisfy the ascending chain condition on right ideals.

    Therefore $R$ is not right Noetherian.
\end{proof}

\begin{lemma}
    Let $R$ be a left Artinian ring. Then its Jacobson radical $J(R)$ is nilpotent.
\end{lemma}

\begin{proof}
    Set
    \[
        J:=J(R).
    \]
    Since $R$ is left Artinian, the descending chain of left ideals
    \[
        J \supseteq J^2 \supseteq J^3 \supseteq \cdots
    \]
    stabilizes. Hence there exists some $n \ge 1$ such that
    \[
        J^n = J^{n+1}.
    \]
    Let
    \[
        N:=J^n.
    \]
    We claim that $N=0$.

    Suppose, for contradiction, that $N \neq 0$. Then
    \[
        N^2 = J^{2n} = J^n = N,
    \]
    since the chain has stabilized.

    Now consider the set
    \[
        \Sigma=\{\, I \leq {}_R R \mid NI \neq 0 \,\}
    \]
    of left ideals of $R$. This set is nonempty because
    \[
        N^2 = N \neq 0,
    \]
    so $N \in \Sigma$.

    Since $R$ is left Artinian, the set of left ideals satisfies the descending chain
    condition, so $\Sigma$ contains a minimal element. Let $I_0 \in \Sigma$ be such a
    minimal left ideal.

    Because $NI_0 \neq 0$, there exists some $a \in I_0$ such that
    \[
        Na \neq 0.
    \]
    Then $Na$ is a left ideal of $R$ contained in $I_0$, and moreover
    \[
        N(Na)=N^2a=Na\neq 0.
    \]
    Thus $Na \in \Sigma$. By the minimality of $I_0$, we must have
    \[
        Na=I_0.
    \]

    Since $a \in I_0 = Na$, there exists $b \in N$ such that
    \[
        a=ba.
    \]
    Therefore
    \[
        (1-b)a=0.
    \]
    But $b \in N \subseteq J(R)$, and for every element of the Jacobson radical,
    $1-b$ has a left inverse. Hence there exists $c \in R$ such that
    \[
        c(1-b)=1.
    \]
    Multiplying the equation $(1-b)a=0$ on the left by $c$, we get
    \[
        a=c(1-b)a=0,
    \]
    contradicting the choice of $a$ with $Na \neq 0$.

    This contradiction shows that $N=0$. Hence
    \[
        J^n=0,
    \]
    so $J(R)$ is nilpotent.
\end{proof}

\begin{theorem}[Levitzki]
    Every left Artinian ring is left Noetherian.
\end{theorem}

\begin{proof}
    Let $R$ be a left Artinian ring, and let $J=J(R)$ be its Jacobson radical.

    Since $R$ is left Artinian, the quotient ring $R/J$ is also left Artinian. Moreover,
    $J(R/J)=0$, hence $R/J$ is semisimple. Therefore every left $R/J$-module is
    semisimple.

    Also, since $R$ is left Artinian, $J$ is nilpotent. Thus there exists $n\ge 1$ such
    that
    \[
        J^n=0.
    \]
    Consider the finite chain of left ideals
    \[
        0=J^n \subseteq J^{n-1} \subseteq \cdots \subseteq J \subseteq R.
    \]

    For each $i=0,1,\dots,n-1$, the factor module $J^i/J^{i+1}$ is naturally a left
    $R/J$-module, because
    \[
        J\cdot J^i \subseteq J^{i+1}.
    \]
    Hence $J^i/J^{i+1}$ is semisimple.

    On the other hand, $J^i/J^{i+1}$ is Artinian, being a quotient of the Artinian
    left $R$-module ${}_R R$. A semisimple Artinian module must be a finite direct sum
    of simple modules, so it is Noetherian. Therefore each factor
    \[
        J^i/J^{i+1}
    \]
    is Noetherian.

    Now we argue by descending induction on $i$. Since $J^n=0$ is Noetherian, and for
    each $i=0,1,\dots,n-1$ we have a short exact sequence
    \[
        0 \longrightarrow J^{i+1} \longrightarrow J^i \longrightarrow J^i/J^{i+1}
        \longrightarrow 0,
    \]
    it follows that if $J^{i+1}$ is Noetherian, then so is $J^i$. Hence
    \[
        J^{n-1},\,J^{n-2},\,\dots,\,J,\,R
    \]
    are all Noetherian as left $R$-modules.

    In particular, ${}_R R$ is Noetherian. Therefore $R$ is a left Noetherian ring.
\end{proof}

\begin{theorem}
    If $R$ is left Noetherian/Artinian, then every finitely generated left $R$-module is Noetherian/Artinian.
\end{theorem}

\begin{proof}
    Prove for Noetherian case: If $M$ is generated by $m$ elements, then $M=Rm_1+ \ldots + Rm_n$, each $Rm_i$ is a homomorphic image of ${}_RR$. Hence $Rm_i$ is Noetherian.
    Finite sums of Noetherian modules are Noetherian, so $M$ is Noetherian.
\end{proof}

